% concatenated from main.tex, extra-package.tex, abstract-v1.tex, intro-v1.tex, recap-A.tex, model-v1.tex, main-pre.tex, section-MR-A.tex, section-MR-B.tex, section-MR-C.tex, conclusion-v1.tex, Appendix-A.tex, Appendix-B.tex, recap-B.tex
\documentclass[conference]{IEEEtran}
\IEEEoverridecommandlockouts
% The preceding line is only needed to identify funding in the first footnote. If that is unneeded, please comment it out.
\newcommand{\tsp}
{\mathsf{\scriptstyle{T}}}
\usepackage{verbatim} % for begin comment
\usepackage{balance}
%---------省略符號------------
\newcommand{\ra}{\rightarrow}
\newcommand{\vs}{\ |\ }
%一般集合指令
\newcommand{\N}{\mathbb{N}}
\newcommand{\Z}{\mathbb{Z}}
\newcommand{\Q}{\mathbb{Q}}
\newcommand{\R}{\mathbb{R}}
\newcommand{\C}{\mathbb{C}}
\newcommand{\M}{\mathbb{M}}
\newcommand{\F}{\mathbb{F}}
\newcommand{\I}{\mathbb{I}}
\newcommand{\E}{\mathbb{E}}
%---------定理宣告------------
\usepackage{amsthm} % theroemstyle 需要使用的套件

%\newtheorem{axm}{Axiom}
\theoremstyle{definition}
\newtheorem{lma}{Lemma}

\theoremstyle{plain}
\newtheorem{thm}{Theorem}
%\numberwithin{thm}{section}
%---------符號宣告------------
\newcommand{\norm}[1]{\left\lVert#1\right\rVert}
%-------------------------------

%\usepackage[utf8]{inputenc}
%\usepackage{fourier} 
\usepackage{array}
\usepackage{makecell}

\renewcommand\theadalign{bc}
\renewcommand\theadfont{\bfseries}
\renewcommand\theadgape{\Gape[4pt]}
\renewcommand\cellgape{\Gape[4pt]}

\renewcommand{\cellalign}{lc}
\renewcommand{\theadalign}{lc}
%%
%% For Wide Tables
\usepackage{blindtext}
%% For Wide Tables
\usepackage{xcolor}
\usepackage{cite}
\usepackage{amsmath,amssymb,amsfonts}
\usepackage{algorithmic}
\usepackage{graphicx}
\usepackage{textcomp}
\usepackage{xcolor,soul}
\def\BibTeX{{\rm B\kern-.05em{\sc i\kern-.025em b}\kern-.08em
    T\kern-.1667em\lower.7ex\hbox{E}\kern-.125emX}}
    
\newcommand{\mh}[1]{\textcolor{red}{#1}}
    
\begin{document}

\title{On the Capacity of Zero-Drift First Arrival Position Channels in Diffusive Molecular Communication
\iffalse
{\footnotesize \textsuperscript{*}Note: Sub-titles are not captured in Xplore and
should not be used}
\thanks{Identify applicable funding agency here. If none, delete this.}
\fi
}

\author{\IEEEauthorblockN{Yen-Chi Lee}
\IEEEauthorblockA{\textit{Hon Hai (Foxconn) Research Institute} \\
%\textit{name of organization (of Aff.)}\\
Taipei, Taiwan \\
%yenchilee1925@gmail.com}
yen-chi.lee@foxconn.com}
\iffalse
\and
\IEEEauthorblockN{ 
Jen-Ming Wu
}
\IEEEauthorblockA{\textit{Hon Hai (Foxconn) Research Institute} \\
%\textit{name of organization (of Aff.)}\\
Taipei, Taiwan \\
jen-ming.wu@foxconn.com}
\fi
\and
\IEEEauthorblockN{Min-Hsiu Hsieh}
\IEEEauthorblockA{\textit{Hon Hai (Foxconn) Research Institute} \\
%\textit{name of organization (of Aff.)}\\
Taipei, Taiwan \\
min-hsiu.hsieh@foxconn.com}
%\thanks{The full version of this paper, including appendices, can be found on arXiv: https://arxiv.org/abs/2201.11383v3. It is referenced in this manuscript as \cite{lee2022capacity}.}
}

\maketitle

\begin{abstract}
%%--

\color{black}
Recent advancements in understanding the impulse response of the first arrival position (FAP) channel in molecular communication (MC) have illuminated its Shannon capacity. While Lee et al. shed light on FAP channel capacities with vertical drifts, the zero-drift scenario remains a conundrum, primarily due to the challenges associated with the heavy-tailed Cauchy distributions whose first and second moments do not exist, rendering traditional mutual information constraints ineffective. This paper unveils a novel characterization of the zero-drift FAP channel capacity for both 2D and 3D. Interestingly, our results reveal a 3D FAP channel capacity that is double its 2D counterpart, hinting at a capacity increase with spatial dimension growth. Furthermore, our approach, which incorporates a modified logarithmic constraint and an output signal constraint, offers a simplified and more intuitive formula (similar to the well-known Gaussian case) for estimating FAP channel capacity.
\color{black}

% Below is longer version provided by GPT4
\iffalse
\color{blue}
Recently, the impulse response of the first arrival position (FAP) channel in molecular communication (MC) has been understood better through new studies. Yet, the Shannon capacity of these FAP channels remains a topic of interest. While Lee et al. have advanced our understanding of the FAP channel capacity, especially in cases with vertical drift toward the receiver, the channel capacity for zero-drift scenarios remains a challenging question.
The core challenge in determining the capacity of zero-drift FAP channels stems from the heavy-tailed nature of the Cauchy distributions. Traditional methods that constrain the mutual information are ineffective because both the first and second moments of the Cauchy distributions are non-existent.
In this paper, we introduce a novel characterization of the zero-drift FAP channel capacity in both 2D and 3D scenarios. Interestingly, our findings suggest that the capacity formula for FAP channels mirrors the Gaussian channel when a second-moment power constraint is applied. Additionally, the capacity of a 3D FAP channel is observed to be double that of its 2D counterpart, reinforcing the idea that FAP channel capacities grow with increasing spatial dimensions.
Key contributions of this work involve the use of a modified logarithmic constraint as an alternative to the conventional power constraint and the adoption of an output signal constraint in place of the traditional input signal constraint. This results in a more streamlined formula for channel capacity.
\color{black}
\fi
%%--


\iffalse
\color{black}
The impulse response of the first arrival position (FAP) channel in molecular communication (MC) has been derived in recent studies. But the Shannon capacity of FAP channels still remains illusive. 
A recent work done by Lee et al. has made strides in characterizing the FAP channel capacity in cases that involves a vertical drift towards the receiver. However, for scenarios where the drift vanishes, the channel capacity problem remains unsolved.
The fundamental obstacle to determining the capacity of FAP channels for zero-drift case is rooted in the heavy-tail nature
concerning the Cauchy distributions. Conventional approaches for constraining the mutual information are inapplicable as both the first and second moments of Cauchy distributions do not existent.
This paper presents a novel characterization of the zero-drift FAP channel capacity for 2D and 3D cases. The capacity formula for the FAP channel is found to have a form similar to the Gaussian channel case under second-moment power constraint. Notably, the capacity of the 3D FAP channel is twice that of the 2D FAP channel, providing evidence that FAP channels have greater capacity as spatial dimensions increase.
Our technical contributions include the application of a modified logarithmic constraint in lieu of the typical power constraint, and the selection of an output signal constraint as a replacement for the input signal constraint, resulting in a more concise formula.
\color{black}
\fi

\end{abstract}

\begin{IEEEkeywords}
diffusive molecular communication, absorbing receivers, first arrival position (FAP), channel capacity, Cauchy distribution.
\end{IEEEkeywords}

\section{Introduction}
\color{black}
Molecular communication (MC) represents a novel communication paradigm rooted in the transmission of information through molecular exchange \cite{nakano2013molecular,yeh2012new}. Owing to its nano-scale applicability and biocompatibility, MC emerges as a potent communication technique for nano-networks \cite{akyildiz2008nanonetworks,farsad2016comprehensive}. Within these systems, tiny message molecules (MM) act as the principal information carriers. A propagation mechanism is necessary for transporting MMs to the receiver (Rx), and this mechanism can be either diffusive (a.k.a. diffusion-based) \cite{pierobon2012capacity}, flow-based \cite{kadloor2012molecular}, or due to an engineered transport system like molecular motors \cite{moore2006design,gregori2010new}. Among these different propagation mechanisms, diffusive MC (sometimes in combination with a drift field) has been the most prevalent approach for MC theoretical research \cite{jamali2019channel}.
\color{black}

\begin{comment}
Molecular communication (MC) is a communication paradigm based on the exchange of molecules \cite{nakano2013molecular,yeh2012new}. Due to the nano-scale feasibility and bio-compatibility, MC is a promising communication approach for nano-networks \cite{akyildiz2008nanonetworks,farsad2016comprehensive}. In MC systems, tiny message molecules (MM) operate as information carriers. A propagation mechanism is necessary for transporting MMs to the receiver (Rx), and this mechanism can be either diffusive (a.k.a. diffusion-based) \cite{pierobon2012capacity}, flow-based \cite{kadloor2012molecular}, or an engineered transport system like molecular motors \cite{moore2006design,gregori2010new}. Among these different propagation mechanisms, diffusive MC, sometimes in combination with a drift field, has been the most prevalent approach for MC theoretical research. 
\end{comment}

\color{black}
The molecular communication (MC) receiver's reception mechanism can be bifurcated into two primary categories, as delineated by \cite{jamali2019channel}: i) passive reception and ii) active reception. Within the realm of active reception, this study considers the widely recognized fully-absorbing Rx \cite{yilmaz2014three}. It is postulated that the Rx possesses the capability to precisely ascertain both the time \cite{srinivas2012molecular} and the position \cite{lee2016distribution,pandey2018molecular,akdeniz2018molecular} at the initial contact of the MM with the Rx. To investigate the Shannon capacity \cite{shannon1948mathematical} for this novel category of position channels, we adopt a straightforward geometric configuration for the receiver, specifically, a spacious receiving plane \cite{andrews2009accurate}. It is worth noting that this spacious plane model can serve as an approximation, particularly when the transmission distance is significantly shorter than the dimensions of the receiver, as discussed in prior works \cite{lee2016distribution,lee2023char}. Note that MMs are removed upon their interaction with the absorptive receiving plane \cite{yilmaz2014three,pandey2018molecular}. For a visual representation of this system model, please refer to Fig.~1.
\color{black}

\begin{comment}
The reception mechanism of a MC receiver can be categorized, according to \cite{jamali2019channel}, into two classes: i) passive reception, and ii) active reception. We consider a common type of active reception called the fully-absorbing Rx \cite{yilmaz2014three}, and assume that the Rx has the ability to measure the time \cite{srinivas2012molecular} and the position \cite{lee2016distribution,pandey2018molecular,akdeniz2018molecular} precisely once the MM first reaches the Rx.
To investigate the Shannon capacity \cite{shannon1948mathematical} for this novel category of position channels, we adopt a straightforward geometric configuration for the receiver, specifically, a spacious receiving plane \cite{andrews2009accurate}. It is worth noting that this spacious plane model can serve as an approximation, particularly when the transmission distance is significantly shorter than the dimensions of the receiver, as discussed in prior works \cite{lee2016distribution,lee2023char}. 
The MMs will be removed once they hit the (absorbing) receiving plane \cite{yilmaz2014three,pandey2018molecular}.
For a visual representation of this concept, please refer to Fig.~1.
\end{comment}


There are various physical attributes of MMs that
can convey information.
The concept of the First Arrival Position (FAP) as an information-carrying property in MC literature was first introduced in the paper by Lee et al. in 2016 \cite{lee2016distribution}. Subsequently, this idea was further developed, particularly in the context of two-dimensional (2D) scenarios, in a later study conducted by Pandey et al. in 2018 \cite{pandey2018molecular}.
\color{black}
While the majority of studies within the realm of MC primarily concentrate on the First Arrival Time (FAT)\footnote{Readers interested in gaining a more comprehensive understanding of FAT-type modulation and its associated channel characteristics are encouraged to refer to \cite{srinivas2012molecular,li2014capacity} for detailed information.}
for absorbing receivers, it is noteworthy to acknowledge at least two reasons for giving due consideration to FAP as a viable alternative.
\begin{itemize}
    \item[1)] For every independent channel use—specifically, for each transmission of a single MM—the FAT information is one-dimensional (1D). In contrast, the FAP-type modulation can carry information in higher dimensions, up to \(n-1\), in an \(n\)-dimensional space. Consequently, the capacity of FAP channels, for each individual channel use, can increase significantly in higher-dimensional scenarios. This potentially allows FAP channels to surpass FAT channels in terms of capacity. We will establish the capacity formulas for FAP channels in Section~\ref{sec:MR} and validate this assertion for 2D and 3D cases.
    \item[2)] The second consideration pertains to time efficiency. In scenarios involving multiple-MM transmission, MMs may arrive out of sequence due to the inherent unpredictability of the diffusion process, leading to cross-over effects as documented in \cite{eckford2008molecular,jamali2019channel}. To mitigate cross-over effects in the MC system, the transmission duration for each symbol must not be overly short.\footnote{A simple way to conceptualize this is to imagine the necessity for a \emph{guard interval} between successive timing symbols.} Given this, for applications where time efficiency is paramount, exploring FAP-type modulation as an alternative solution might be advisable.
\end{itemize}

\color{black}
In this paper, our primary attention centers on one-shot transmission \cite{murin2018optimal}, referring specifically to transmission using a singular MM. To delve into the channel capacity of either FAP or FAT channels, initiating with a quantitative delineation of the channel impulse response is imperative. The one-shot FAT channel can be succinctly characterized as a time-invariant additive channel \cite{srinivas2012molecular}:
$
    t_{\text{out}} = t_{\text{in}} + t_{\text{n}},
%\label{eq:add-time-ch}
$
where $t_{\text{in}}$ represents the release time, $t_{\text{out}}$ denotes the arrival time, and $t_{\text{n}}$ is the random time delay resulting from propagation dynamics. As established in \cite{srinivas2012molecular}, the variable $t_{\text{n}}$ adheres to the inverse Gaussian distribution. This channel is thus recognized as the Additive Inverse Gaussian Noise (AIGN) channel within the MC field. Subsequent research in \cite{srinivas2012molecular,li2014capacity} has presented bounds on the capacity of the AIGN channel and has further characterized the capacity-achieving input time distribution.
\color{black}

\begin{comment}
In the remaining of this paper, we will mainly focus on one-shot transmission \cite{murin2018optimal} (i.e. transmission using a single MM).
In order to explore the channel capacity of FAP or FAT channels, the first step is to provide quantitative descriptions about the channel impulse response.
The one-shot FAT channel can be described as a time-invariant additive channel \cite{srinivas2012molecular}:
$
    t_{\text{out}} = t_{\text{in}} + t_{\text{n}},
%\label{eq:add-time-ch}
$
where $t_{\text{in}}$ is the releasing time, $t_{\text{out}}$ is the arriving time, and $t_{\text{n}}$ is the random time delay due to the propagation mechanisms. The authors of \cite{srinivas2012molecular} showed that $t_{\text{n}}$ follows the inverse Gaussian distribution. Hence, this channel
is known as the Additive Inverse Gaussian Noise (AIGN) channel in MC literature.
Later in \cite{srinivas2012molecular,li2014capacity}, some bounds on capacity of AIGN channel is derived, and the capacity-achieving input time distribution was also characterized.
\end{comment}

%%----------

\color{black}
For FAP channels operating within \(n\)-dimensional spaces, the one-shot channel model is elegantly defined through an additive vector form, as detailed in \cite{lee2022arrival}:
$
    \mathbf{x}_{\text{out}} = \mathbf{x}_{\text{in}} + \mathbf{x}_{\text{n}},
$
where \(\mathbf{x}_{\text{in}}\) signifies the releasing position, \(\mathbf{x}_{\text{out}}\) indicates the arriving position, and \(\mathbf{x}_{\text{n}}\) captures the random position deviation induced by propagation mechanisms. 
It is crucial to recognize that the vectors \(\mathbf{x}_{\text{in}}\), \(\mathbf{x}_{\text{out}}\), and \(\mathbf{x}_{\text{n}}\) reside within the Euclidean vector space \(\mathbb{R}^{n-1}\). This configuration arises since one dimension is earmarked for physical movement, effectively excluding it from bearing information in that particular dimension.
Though past studies have articulated the probability density function of \(\mathbf{x}_{\text{n}}\) \cite{lee2016distribution,pandey2018molecular,lee2022arrival}, the quest to comprehend the channel capacity of this distinctive position channel is ongoing. A recent exploration by Lee et al. \cite{lee2023char} has elucidated the FAP channel capacity in contexts characterized by a vertical drift toward Rx. However, the channel capacity remains elusive in zero-drift scenarios. This manuscript aims to address this lacuna, unveiling precise formulas for the zero-drift FAP channel capacity for 2D and 3D dimensions, and introducing a novel logarithmic constraint, as detailed in Sections~III and IV.
\color{black}

\begin{comment}
For FAP channels in \(n\)-dimensional spaces, the one-shot channel model can be cast in an additive vector form, as articulated in \cite{lee2022arrival}:
$
    \mathbf{x}_{\text{out}} = \mathbf{x}_{\text{in}} + \mathbf{x}_{\text{n}},
$
where \(\mathbf{x}_{\text{in}}\) is the releasing position, \(\mathbf{x}_{\text{out}}\) denotes the arriving position, and \(\mathbf{x}_{\text{n}}\) encapsulates the random position bias due to propagation dynamics. 
It's pivotal to understand that vectors \(\mathbf{x}_{\text{in}}\), \(\mathbf{x}_{\text{out}}\), and \(\mathbf{x}_{\text{n}}\) are members of the Euclidean vector space \(\mathbb{R}^{n-1}\); this is because one dimension is allocated for traversal, rendering it void of information-carrying capacity in that specific dimension.
While prior research has derived the probability density function of \(\mathbf{x}_{\text{n}}\) \cite{lee2016distribution,pandey2018molecular,lee2022arrival}, discerning the channel capacity for this particular position channel remains an active area of inquiry. A recent contribution by Lee et al. \cite{lee2023char} offers insights into the FAP channel capacity under positive vertical drift conditions. Yet, the channel capacity remains illusive in zero-drift scenarios. This manuscript endeavors to bridge gap, presenting definitive formulas for the zero-drift FAP channel capacity in 2D and 3D contexts, integrating a novel logarithmic constraint, as expounded in Sections~III and IV.
\end{comment}

\begin{comment}
Regarding FAP channels in $n$-dimensional spaces, it is noteworthy that the one-shot channel model can also be expressed in an additive vector form, as elaborated in \cite{lee2022arrival}, which can be succinctly represented as follows:
$
    \mathbf{x}_{\text{out}} = \mathbf{x}_{\text{in}} + \mathbf{x}_{\text{n}},
$
Herein, $\mathbf{x}_{\text{in}}$ denotes the releasing position, $\mathbf{x}_{\text{out}}$ represents the arriving position, and $\mathbf{x}_{\text{n}}$ signifies the random position bias arising from propagation mechanisms. It is essential to emphasize that each of these vectors, $\mathbf{x}_{\text{in}}$, $\mathbf{x}_{\text{out}}$, and $\mathbf{x}_{\text{n}}$, is an element of the Euclidean vector space $\mathbb{R}^{n-1}$.


While the probability density function of $\mathbf{x}_{\text{n}}$ has been derived in previous works \cite{lee2016distribution,pandey2018molecular,lee2022arrival}, the determination of the channel capacity for this specific type of position channel remains a subject of ongoing investigation. Notably, a recent study by Lee et al. \cite{lee2023char} has made strides in characterizing the FAP channel capacity in cases involving positive vertical drift. However, for scenarios where the drift is zero, the channel capacity problem remains unresolved.
This paper is dedicated to addressing this knowledge gap by furnishing explicit formulas for the zero-drift FAP channel capacity in both 2D and 3D spatial dimensions, incorporating a novel logarithmic constraint, as delineated in Sections~III and IV.
\end{comment}
%%----------


%\color{blue}
The remainder of this paper is structured as follows. 
Section~\ref{pre-A} provides a succinct review of mutual information and the Shannon channel capacity framework. Section~\ref{sec:CM} outlines the channel model under scrutiny and introduces the \(\alpha\)-power constraint.
Our core findings are presented in Section~\ref{sec:MR}. We draw our conclusions in Section~\ref{sec:conclude}.
\color{black}

%\mh{The remainder of this paper is structured as follows. (to be added)}

\section{Preliminary} \label{pre-A}
\color{black}
This section revisits foundational concepts surrounding mutual information and channel capacity. Consider \(X\) as a random variable with corresponding probability density function (pdf) \(f\). The support of \(f\) is represented as \(\mathcal{X}\). The differential entropy \(h(X)\) can be expressed as
$
h(X) := \E_X[-\ln f(X)] = -\int_{\mathcal{X}}  [\ln f(x)] f(x)dx.
$
Leveraging the above entropy definition, the mutual information between the input and output of the additive channel \(Y = X + N\) is:
\begin{align}
\begin{split}
    I(X;Y) &= h(Y) - h(Y|X)
    = h(Y) - h(X+N|X) \\
    &= h(Y) - h(X|X) - h(N|X) \\
    &= h(Y) - h(N).
\end{split}
\end{align}
Here, the assumption is that noise \(N\) operates independently of signal \(X\).
Remember that the (Shannon) channel capacity is captured as
$
C = \sup_{f_X(x)} I(X;Y).
\label{eq:inpc}
$
The supremum is evaluated over all possible \(X\) distributions under a selected constraint. For point-to-point communication or single-input-single-output (SISO) scenarios, output constraints can be adopted interchangeably, resulting in
$
C = \sup_{f_Y(x)} I(X;Y).
\label{eq:outpc}
$

Constraint selection is intimately tied to the noise distribution. To enhance understanding, we'll first contemplate a straightforward instance prior to diving into the specifics of FAP channels. For the additive Gaussian channel, where noise follows a Gaussian distribution, the prevailing power measure embraces the second moment constraint:
$
    \E[X^2] = \int_{\mathcal{X}} x^2 f(x)dx \leq P.
$
The constant \(P \geq 0\) denotes the upper power limit for all input signals. Based on this constraint, the channel capacity can be expressed as:
$
C = \sup_{f_X(x) :\  \E[X^2]\leq P} I(X;Y).
$
Consulting a standard reference such as \cite{cover1999elements}, the supremum of \(I(X;Y)\) is unveiled as
$
\frac{1}{2}\ln\left(\frac{\sigma^2+P}{\sigma^2}\right),
$
with the capacity-achieving distribution being Gaussian with variance \(P\), i.e. \(X \sim \mathcal{N}(0,P)\).

For subsequent comparison (presented in Section~\ref{sec:MR}) purpose, we can set \(A^2:=\sigma^2+P\), yielding
\begin{align}
\begin{split}
C_{\text{G}}
=\frac{1}{2}\ln\left(\dfrac{A^2}{\sigma^2}\right)
=\ln\left(\dfrac{A}{\sigma}\right),
\end{split}
\end{align}
where the subscript \(\text{G}\) hints the Gaussian channel.
\color{black}


\begin{comment}
We briefly review some background knowledge for mutual information and channel capacity.
Let $X$ be a random variable with a probability density function $f$ whose support is denoted by $\mathcal{X}$.
The (differential) entropy $h(X)$ is defined as: 
\begin{equation}
h(X):=\E_X[-\ln f(X)]=-\int_{\mathcal{X}}  [\ln f(x)] f(x)dx.
\end{equation}
Based on the definition of entropy, the mutual information between the input and output of the additive channel 
$
    Y=X+N
$
is given by
\begin{align}
\begin{split}
    I(X;Y) &= h(Y) - h(Y|X)
    =h(Y) - h(X+N|X)\\
    &=h(Y) - h(X|X) - h(N|X)\\
    &=h(Y) - h(N).
\end{split}
\end{align}
Here we have assumed that the noise $N$ is independent of signal $X$.
Recall that the (Shannon) channel capacity is defined as:
$
C = \sup_{f_X(x)} I(X;Y).
\label{eq:inpc}
$
The supremum is taken over all possible distributions of $X$ satisfying some chosen constraint. For point-to-point communication or single-input-single-output (SISO) scenarios, sometimes the output constraints are used instead of the input constraints, yielding
$
C = \sup_{f_Y(x)} I(X;Y).
\label{eq:outpc}
$
%%--------------------

The choice of constraint heavily relies on the noise distribution. To gain a better understanding of this concept, we will examine a simple example before delving into the specifics of FAP channels. For additive Gaussian channel (i.e. the noise is Gaussian distributed), the commonly employed power measure is the second moment constraint. We can prescribe our signal set to satisfy:
$
    \E[X^2] 
    = \int_{\mathcal{X}} x^2 f(x)dx \leq P.
%\label{eq:2pc}
$
The constant $P\geq 0$ is the maximum power we allowed for all input signals.
Using this power constraint, the channel capacity can be written as:
$
C = \sup_{f_X(x) :\  \E[X^2]\leq P} I(X;Y).
$
Following a standard textbook derivation (see \cite{cover1999elements}), one can see that the supremum value of $I(X;Y)$ is
$
\frac{1}{2}\ln\left(\frac{\sigma^2+P}{\sigma^2}\right), 
%\label{eq:13}
$
and the capacity attaining distribution is normal distributed with variance $P$, namely, $X \sim \mathcal{N}(0,P)$.


For later comparison (presented in Section~\ref{sec:MR}) to the capacity formula of zero-drift FAP channels, we can let $A^2:=\sigma^2+P$, yielding
\begin{align}
\begin{split}
C_{\text{G}}
=C(A,\sigma)
=\frac{1}{2}\ln\left(\dfrac{A^2}{\sigma^2}\right)
=\ln\left(\dfrac{A}{\sigma}\right),
\end{split}
\end{align}
where the subscript $\text{G}$ stands for Gaussian channel.
\end{comment}


\section{Channel Model} \label{sec:CM}
\begin{figure}[!t]  %[htbp]
\centerline{\includegraphics[width=0.48\textwidth]{./Figures/new2Da.eps}}
\caption{
This figure presents an abstract FAP channel model featuring hyperplane-shaped transmitters (Tx) and receivers (Rx). In this model, the transmitter is assumed to be transparent, permitting particles to traverse it without encountering any forces. It is important to observe that the emission point resides on the Tx hyperplane, and the Rx registers the arrival point when the information particle initially hits the Rx. Note that we consider \emph{zero drift} case for our main results.
}
\label{fig:2Da}
\end{figure}
%%--

We consider an additive vector channel:
$
    \mathbf{Y}=\mathbf{X}+\mathbf{N}
%\label{eq:17}
$
where the noise $\mathbf{N}$ is (multivariate) Cauchy distributed.\footnote{
We postpone the reason why FAP channels reduce to Cauchy channels under zero-drift assumption to Section~\ref{subsec:MR-A}.
}
Because the Cauchy distributions\footnote{
For a detailed characterization of Cauchy density functions, one can refer to Appendix~A.
} are heavy-tailed \cite{sigman1999primer}, traditional signal processing theory, which is tailored for finite second moment signals, does not apply directly to the Cauchy case.

Of the heavy-tailed probability models, the alpha-stable family \cite{pierce1997application} stands out, having demonstrated its prowess in simulating multiple access interference \cite{el2010alpha} and co-channel interference from a Poisson-distributed field of interferers \cite{gulati2010statistics}. In alpha-stable channels, the second moment is not a viable power measure due to its inherent infinity.

\color{black}
The task of discerning the channel capacity for alpha-stable additive noise—specifically when \(\alpha \geq 1\)—under an \(r\)-th absolute moment input constraint was adeptly addressed by \cite{fahs2012capacity} for symmetric alpha-stable noise scenarios. Their exploration led to the identification of an optimal capacity-achieving input that is both discrete and compactly defined. Meanwhile, a distinct concept of power, termed as the ``geometric power," was ushered in by \cite{gonzalez1997zero} to cater to heavy-tailed distributions \cite{sigman1999primer}. Within the context of this manuscript, we introduce an alternative power characterization schema for the \(\alpha\)-stable family, which is referred to as the \(\alpha\)-power, as detailed in \cite{fahs2017information}. It's pertinent to highlight that while the genesis of this framework traces back to \cite{fahs2017information}, our work repurposes its essence to adeptly navigate the challenges posed by zero-drift FAP channels in the realm of molecular communication.
\color{black}

\begin{comment}
The problem of finding the channel capacity for alpha-stable additive noise (with $\alpha \geq 1$) with an $r$-th absolute moment input constraint was solved by \cite{fahs2012capacity} for the case of symmetric alpha-stable noise, and the optimal capacity-achieving input was found to be discrete and compactly supported. An alternative notion of power, called the geometric power, was introduced by \cite{gonzalez1997zero} for heavy-tailed distributions \cite{sigman1999primer}. In this paper, we introduce another power characterization framework for $\alpha$-stable family called the $\alpha$-power \cite{fahs2017information}. 
(Note that this framework is originally proposed in \cite{fahs2017information}, and we adopt this framework to tackle with the zero-drift FAP channels in molecular communication field.)
\end{comment}

%%-----------------------
Let us define a relative power measure $P_{\alpha}(\mathbf{X})$ satisfying the following:
\begin{itemize}
    \item[P1)] $P_\alpha(\mathbf{X}) \geq 0$, with equality if and only if $\mathbf{X}=0$ almost surely.
    \item[P2)] $P_\alpha(k\mathbf{X})=|k|P_\alpha(\mathbf{X})$, for any $k\in\R$.
\end{itemize}
For $\alpha$-stable distributions, it is known (see \cite{fahs2017information}) that
\begin{equation}
\begin{cases}
\E_\mathbf{X}\big[\ln(1+\norm{\mathbf{X}}^2)\big]<\infty,\quad &\alpha<2;\\
\E_\mathbf{X}\big[\norm{\mathbf{X}}^2\big] < \infty,\quad &\alpha=2.
\end{cases}
\end{equation}
We choose the power measure $P_1(\cdot)$ for multivariate Cauchy-like distributions as follows:
\begin{equation}
    \E_\mathbf{X}\left[\ln\left(1+\norm{\frac{\mathbf{X}}{P_1(\mathbf{X})}}^2 \right)\right]
    = w_2\left(\frac{1+p}{2};\frac{p}{2}\right),
\end{equation}
where the constant evaluation function $w_2(t;\alpha)$ can be written explicitly as
the difference of two digamma functions:
$
    w_2(t;a) = \psi(t) - \psi(t-a), \text{\ for\ } t>a.
$
Note that the values of $w_2$ function at $p=1$ and $p=2$ are $2\ln(2)$ and $2\ln(e)$ respectively. 

\color{black}

\color{black}
\section{Main Results}\label{sec:MR}
In this section, we delve into the intricacies of the FAP communication channel. We commence by demonstrating in Section~\ref{subsec:MR-A} that as the drift velocity in the fluid medium approaches zero, the 2D and 3D FAP channels simplify to univariate and bivariate Cauchy distributions, respectively.

Proceeding to Sections~\ref{MR-B} and~\ref{MR-C}, we integrate the concept of the \(\alpha\)-power, a relative power measure documented in \cite{fahs2017information}, into MC discussions. We explicitly define the signal space requisite to depict the no-drift FAP channels. Employing our introduced logarithmic constraint on the output signal space, we ascertain the Shannon capacity of the FAP channel for both 2D and 3D configurations in a closed-form manner. Key insights regarding the FAP channel capacity are encapsulated in Theorem~\ref{thm:T1} and Theorem~\ref{thm:T2}. 

Notice that in the literature, while \cite{fahs2014cauchy} provided a capacity formula for the univariate Cauchy channel (not for MC purpose) under a logarithmic constraint applied to the input signal space, directly extending their findings to higher-dimensional scenarios is challenging due to intricate integration procedures, as evident in \cite[Appendix~I]{fahs2014cauchy}. One way to tackle this is replacing the input constraints with output constraints, as demonstrated in Sections~\ref{MR-B} and~\ref{MR-C}.
\color{black}

\begin{comment}
\section{Main Results}\label{sec:MR}
%\color{blue}
We consider the FAP communication channel. In this main section, we first show in Section~\ref{subsec:MR-A} that when the drift velocity in the fluid medium approaches zero, both the 2D and 3D FAP channel reduce to univariate and bivariate Cauchy distributions respectively. This result is new to the MC society.

Next in Section~\ref{MR-B} and Section~\ref{MR-C}, we introduce a relative power measure called the $\alpha$-power (see \cite{fahs2017information}) into the MC field. We write down explicitly the signal space needed to characterize the no-drift FAP channels. Under this proposed logarithmic constraint on the output signal space, we derive the Shannon capacity for FAP channel both in 2D and 3D spaces in closed-form. The conclusions about the FAP channel capacity are summarized in Theorem~\ref{thm:T1} and Theorem~\ref{thm:T2}.
\color{black}
Notice that in \cite{fahs2014cauchy}, the authors derived the capacity formula for univariate Cauchy channel under logarithmic constraint imposed on the input signal space. However, their result is not easy to extend to high dimensional cases because a complicated integration process was involved, see \cite[Appendix~I]{fahs2014cauchy}.
\color{black}
\end{comment}

\subsection{FAP Channel Reduces to Cauchy Channel Under Zero Drift Condition}\label{subsec:MR-A}
\subsubsection{In 2D Space}
From \cite{lee2022arrival} we know that
the FAP density function in 2D space, with a drift velocity $\mathbf{v}=(v_1,v_2)$, can be written as:
\begin{align}
\begin{split}
f_{Y \mid X}(\mathbf{y}|\mathbf{x})=&\ \dfrac{|\mathbf{v}|\lambda}{\sigma^2\pi}\exp\left\{\dfrac{-v_2 \lambda}{\sigma^2}\right\} \exp\left\{\dfrac{-v_1(x_1-y_1)}{\sigma^2}\right\}\\ &\cdot\dfrac{K_1\left(\dfrac{|\mathbf{v}|}{\sigma^2}\sqrt{(x_1-y_1)^2+\lambda^2}\right)}{\sqrt{(x_1-y_1)^2+\lambda^2}},
\end{split}
\label{eq:2DFAP}
\end{align}
where $\mathbf{x}=(x_1,\lambda)$, $\mathbf{y}=(y_1,0)$, and $\lambda$ is the \textit{transmission distance} between transmitter (Tx) and Rx. Note that in Eq.~\eqref{eq:2DFAP}, $\sigma^2$ is the microscopic diffusion coefficient which can be related to the macroscopic diffusion coefficient $D$ through the relation $\sigma^2=2D$, see \cite{farsad2016comprehensive}.
Note also that the special function $K_{1}(\cdot)$ in Eq.~\eqref{eq:2DFAP} is the modified Bessel function of the second kind (see \cite{bell2004special}) with order $\nu=1$.

Next we show that when drift velocity approaches zero, the 2D FAP density reduces to a univariate Cauchy distribution. 
We use a limit property about the special function $K_1(\cdot)$ from \cite{yang2017approximating}:
$
    \lim_{x\xrightarrow{}0} xK_1(x)=1.
%\label{eq:limit-prop}
$
Using this limit property, we can do the following calculation:
\begin{align}
\begin{split}
&\ f_{Y \mid X; \mathbf{v}=\mathbf{0}}(\mathbf{y}|\mathbf{x})\\
=&\ 
\lim_{|\mathbf{v}|\xrightarrow{}\mathbf{0}}
\frac{|\mathbf{v}|\lambda}{\sigma^2\pi}
\cdot\frac{K_1\left(\frac{|\mathbf{v}|}{\sigma^2}\sqrt{(x_1-y_1)^2+\lambda^2}\right)}{\sqrt{(x_1-y_1)^2+\lambda^2}}\\
=&\ 
\dfrac{\lambda}{\pi}
\lim_{|\mathbf{v}|\xrightarrow{}\mathbf{0}} \Bigg[
\frac{|\mathbf{v}|}{\sigma^2}\sqrt{(x_1-y_1)^2+\lambda^2}\\ &\quad\quad\quad\quad\quad \cdot\frac{K_1\left(\frac{|\mathbf{v}|}{\sigma^2}\sqrt{(x_1-y_1)^2+\lambda^2}\right)}{(x_1-y_1)^2+\lambda^2}
\Bigg]
\\
=&\ \dfrac{\lambda}{\pi}
\frac{1}{(x_1-y_1)^2+\lambda^2}.
\end{split}
\label{eq:2D-zero}
\end{align}
This calculation result can be interpreted as follows. Let $x_1$ be the input position and $y_1$ be the output position. The displacement of position is an additive noise $n$, following the density function:
$
\frac{\lambda}{\pi}
\frac{1}{n^2+\lambda^2},
$
which is Cauchy distributed. Please refer to Appendix~A for the form of density function of Cauchy random variables.

\subsubsection{In 3D Space}
From \cite{lee2022arrival,lee2016distribution} we know that
the FAP density function in 3D diffusion channel, with a drift velocity $\mathbf{v}=(v_1,v_2,v_3)$, can be written as:
\begin{equation}
\begin{aligned}
&f_{Y \mid X}(\mathbf{y}|\mathbf{x})\\
=\ & \frac{\lambda}{2 \pi} \exp \left\{-\dfrac{v_{3} \lambda}{\sigma^2}\right\} \exp \left\{\dfrac{v_{1}}{\sigma^2}\left(y_1-x_1\right)+\dfrac{v_{2}}{\sigma^2}\left(y_2-x_2\right)\right\} \\
& \cdot \exp \left\{-\dfrac{|\mathbf{v}|}{\sigma^2}\norm{\mathbf{y}-\mathbf{x}}\right\}\left[\frac{1+\frac{|\mathbf{v}|}{\sigma^2}\norm{\mathbf{y}-\mathbf{x}}}{\norm{\mathbf{y}-\mathbf{x}}^{3}}\right],
\end{aligned}
\label{eq:3DFAP}
\end{equation}
where $\sigma^2$ is the microscopic diffusion coefficient, the scale parameter $\lambda$ is the \textit{transmission distance} between Tx and Rx plane, $v_3$ is the longitudinal component of drift velocity in the direction parallel to the transmission direction, and $v_1$, $v_2$ are the components perpendicular (or transverse) to the transmission direction. Here we use the symbol $\norm{\cdot}$ to represent the Euclidean norm. Namely,
$
\norm{\mathbf{y}-\mathbf{x}}=\sqrt{(y_1-x_1)^2+(y_2-x_2)^2+\lambda^2},
$
where $\mathbf{x}=(x_1,x_2,\lambda)$, $\mathbf{y}=(y_1,y_2,0)$ are position vectors in $\R^3$.

Next we show that when drift velocity approaches zero, the 3D FAP density reduces to a bivariate Cauchy distribution.
Note that when $\mathbf{v}=\mathbf{0}$, all the exponential terms in Eq.~\eqref{eq:3DFAP} become $e^0=1$,
so that we have:
\begin{align}
\begin{split}
&\ f_{Y \mid X;\mathbf{v}=\mathbf{0}}(\mathbf{y}|\mathbf{x})\\
=&\ 
\frac{\lambda}{2 \pi}
\lim_{|\mathbf{v}|\xrightarrow{}\mathbf{0}}
\frac{1+\frac{|\mathbf{v}|}{\sigma^2}\norm{\mathbf{y}-\mathbf{x}}}{\norm{\mathbf{y}-\mathbf{x}}^{3}}
=\
\frac{\lambda}{2 \pi}
\frac{1}{\norm{\mathbf{y}-\mathbf{x}}^{3}}\\
=&\
\dfrac{\lambda}{2\pi}
\dfrac{1}{
\big((y_1-x_1)^2+(y_2-x_2)^2+\lambda^2\big)^{3/2}
}.
\end{split}
\end{align}
That is, if we regard the position channel as:
$
\mathbf{y}=\mathbf{x}+\mathbf{n},
$
then $\mathbf{n}$ follows a bivariate Cauchy distribution, see Eq.~\eqref{eq:simple-bi-ch} in Appendix~A for its pdf.
\color{black}

%\subsection{Logarithmic Power Measure for Multivariate Cauchy Distributions}

\subsection{Channel Capacity for 2D FAP Channel Under Logarithmic Constraint}\label{MR-B}
%\color{blue}
As mentioned in Section~\ref{subsec:MR-A}, a 2D FAP channel reduces to a Cauchy channel when the drift velocity approaches zero. Since we are considering a point-to-point (or SISO) communication scenario, the input constraint is actually equivalent to the output constraint. We choose to adopt output constraint to simplify the calculation.

It is known that a Cauchy distribution $X \sim {\rm Cauchy}(0,k)$ maximizes the entropy among all RVs\footnote{This is called the \textit{principle of maximum entropy} (see \cite{kapur1989maximum}) for Cauchy distribution.} that satisfy 
\begin{equation}
\E_X \ln \left[1+\left(\frac{X}{k}\right)^2\right] = 2\ln(2),
\end{equation}
or equivalently,
\begin{equation}
\E_X \ln \left[1+X^2\right] = 2\ln(1+k).
\end{equation}
The corresponding maximum entropy value is $\ln(4\pi k)$, see Appendix~A. Note that
the parameter $k$ restricts the dispersion of all feasible random variables $X$.

%Consider the additive channel $Y=X+N$ where $X$, $Y$ are the input and output of the channel respectively, and $N$ is a Cauchy RV, i.e. $N \sim {\rm Cauchy}(n_0,\gamma)$. Notice that a translation of the probability density function does not affect the mutual information, namely $I(X;Y)=I(X;Y-n_0)$, so we can assume without loss of generality that $N \sim {\rm Cauchy}(0,\gamma)$. The main theorem in [1] is restated as follows.

Consider a 2D FAP zero-drift channel (as shown in Section~\ref{subsec:MR-A}), the input-output relation can be written as
\begin{equation}
    Y=X+N,
\label{eq:29}
\end{equation}
where $N\sim {\rm Cauchy}(0,\lambda)$.
For continuous-variable channel capacity problem, we need to specify a family of distributions (i.e. feasible set) that are under consideration in order to prevent the capacity value from being infinite. Instead of writing down explicitly the constraint equality for $X$ or $Y$, we define
\begin{align}
\begin{split}
\mathcal{D}(A)
=\ 
&\Big\{\text{distributions\ }Y\ \Big|\ \exists k\in[\lambda,A]\\ 
&\quad\quad\quad\text{\ such that\ }
\E_Y \ln \Big[1+\Big(\frac{Y}{k}\Big)^2\Big] = 2\ln(2)
\Big\}.
\end{split}
\label{eq:domain-Y}
\end{align}
The parameter $A$ appeared in Eq.~\eqref{eq:domain-Y} indicates the ``largest allowed dispersion'' for the distributions in $\mathcal{D}(A)$. 
For later use, we also define
\begin{align}
\begin{split}
\mathcal{D}_k
=\ 
&\Big\{\text{distributions\ }Y\ \Big|\
\E_Y \ln \Big[1+\Big(\frac{Y}{k}\Big)^2\Big] = 2\ln(2)
\Big\},
\end{split}
\label{eq:domain-k}
\end{align}
so that we have
$
    \mathcal{D}(A)=\bigcup_{k\in[\lambda,A]} \mathcal{D}_k.
$
With these notations, now we can state our main results in the following two theorems.
\begin{thm}
The capacity of 2D FAP channel \eqref{eq:29} is
\begin{equation}
    C_{\rm{2D,FAP}}=C(A,\lambda)=\ln\left(\dfrac{A}{\lambda}\right)
\end{equation}
under the output logarithmic constraint
$
Y\in \mathcal{D}(A)
$
for some 
prescribed dispersion level $A$, where $A \geq \lambda$.
In addition, the corresponding capacity achieving output distribution is
$
Y^* \sim {\rm Cauchy}(0,A),  
$
or equivalently,
$
X^* \sim {\rm Cauchy}(0,A-\lambda).
$
\label{thm:T1}
\end{thm}

\color{black}
Given that the proof of Theorem~\ref{thm:T1} closely parallels that of Theorem~\ref{thm:T2}, we defer the proof details to Appendix~B and proceed directly to the examination of the 3D case.
\color{black}

%Since the proof of Theorem~\ref{thm:T1} is similar to the proof of Theorem~\ref{thm:T2}, we settle the proof details into \cite[Appendix~B]{lee2022capacity} and move directly to the 3D case.

















%%----
\iffalse
A few observations about this theorem are important:
\begin{description}
\item[(i)] Let the function
\begin{equation}
    g(k;x) := \ln \left[
    \Big(\frac{k+\gamma}{k}\Big)^2
    +\Big(\frac{x}{k}\Big)^2
    \right].
\end{equation}
One can show that the value of $g(k;x)$ is decreasing when $k$ is increasing.
\item[(ii)]
We have
\begin{equation}
    g(\gamma;x) := \ln \left[
    4
    +\Big(\frac{x}{\gamma}\Big)^2
    \right] \geq \ln(4).
\end{equation}
\item[(iii)]
Different values of $A$ determines the feasible set of input distributions. When $A<\gamma$, the feasible set prescribed by (21) is the empty set. When $A=\gamma$, the feasible input distribution prescribed by (21) has only one choice: $X=0$ almost surely.
\item[(iv)]
Suppose $A\geq \gamma$ is given.
Each input distribution $X$ satisfying (21) corresponds to a value $k$ such that
\begin{equation}
\E \ln \left[
    \Big(\frac{k+\gamma}{k}\Big)^2
    +\Big(\frac{X}{k}\Big)^2
    \right] = \ln(4),
\end{equation}
where the value of $k$ falls in
the interval $[\gamma,A]$.
\item[(v)]
For any distribution satisfying
(21), the quantity 
\begin{equation}
    P=k-\gamma
\end{equation}
can be viewed as a new type of power measure (we refer to it as \textit{logarithmic power}) of the input distribution. The logarithmic power of these feasible distributions falls in the range $[0,A-\gamma]$.
\end{description}
Here we present an useful lemma before proceeding further.
\begin{lma}
Suppose $X$, $N$ are independent random variables, and $Y=X+N$. The following two expressions are equivalent:
\begin{description}
\item[(a)]
\begin{equation}
\E_X \ln \left[
    \Big(\frac{k+\gamma}{k}\Big)^2
    +\Big(\frac{X}{k}\Big)^2
    \right] = \ln(4)
\end{equation}
\item[(b)]
\begin{equation}
\E_Y \ln \left[
    \Big(
    1
    +\Big(\frac{Y}{k}\Big)^2
    \right] = \ln(4)
\end{equation}
\end{description}
where $k\in[\gamma,A]$.
\label{lma:L4}
\end{lma}
\fi

\subsection{Channel Capacity for 3D FAP Channel Under Logarithmic Constraint}\label{MR-C}
\color{black}
\iffalse
As mentioned in Section~\ref{MR-A}, a 3D FAP channel reduces to bivariate Cauchy channel when the drift velocity is zero. For the problem of finding channel capacity of zero-drift 3D FAP channel, most procedures turn out to be very similar as presented in \ref{MR-C}, except some variations in the constants depending on the spatial dimension. 
\fi

%Our strategy in this subsection is that, we state the different points, and then state the main theorem. The calculation details can be easily reconstructed by the readers with the help of Section~\ref{MR-C}.

It is known (see \cite{kapur1989maximum},\cite{kotz2004multivariate}) that a bivariate Cauchy distribution $\mathbf{X}\sim{\rm Cauchy}_2 (0,\Sigma=\text{diag}(k^2,k^2))$ maximizes the entropy among all bivariate RVs that satisfy 
\begin{equation}
\E_{\mathbf{X}} \ln \left[1+\norm{\frac{\mathbf{X}}{k}}^2\right] = 2\ln\left(e\right),
\end{equation}
The corresponding maximum entropy value is $\ln(2\pi e^3 k^2)$, see Appendix~A. Note that
the parameter $k$ restricts the dispersion of all feasible random vectors $\mathbf{X}$.

Consider a 3D FAP zero-drift channel (as shown in Section~\ref{subsec:MR-A}), the input-output relation can be written as
\begin{equation}
    \mathbf{Y}=\mathbf{X}+\mathbf{N},
\label{eq:45}
\end{equation}
where
$\mathbf{N}\sim{\rm Cauchy}_2 (\mathbf{0},\Sigma=\text{diag}(\lambda^2,\lambda^2))$.
We define
\begin{align}
\begin{split}
\mathcal{D}_{\text{bi}}(A)
=\ 
&\Bigg\{\text{distributions\ }\mathbf{Y}\ \Bigg|\ \exists k\in[\lambda,A]\\ 
&\quad\quad\quad\text{\ such that\ }
\E_\mathbf{Y} \ln \Big[1+\norm{\frac{\mathbf{Y}}{k}}^2\Big] = 2\ln(e)
\Bigg\}.
\end{split}
\label{eq:domain-Y-bi}
\end{align}
The parameter $A$ appeared in Eq.~\eqref{eq:domain-Y-bi} indicates the ``largest allowed dispersion'' for the distributions in $\mathcal{D}_{\text{bi}}(A)$. 
For later use, we also define
\begin{align}
\begin{split}
\mathcal{D}_{k,\text{bi}}
=\ 
&\Bigg\{\text{distributions\ }\mathbf{Y}\ \Bigg|\
\E_\mathbf{Y} \ln \Big[1+\norm{\frac{\mathbf{Y}}{k}}^2\Big] = 2\ln(e)
\Bigg\},
\end{split}
\label{eq:domain-k-bi}
\end{align}
so that we have
\begin{equation}
    \mathcal{D}_{\text{bi}}(A)=\bigcup_{k\in[\lambda,A]} \mathcal{D}_{k,\text{bi}}.
\end{equation}
The 3D FAP capacity formula is stated in the following theorem.
\begin{thm}
The capacity of 3D FAP channel \eqref{eq:45}  is
\begin{equation}
    C_{\rm{3D,FAP}}=C(A,\lambda)=2\ln\left(\dfrac{A}{\lambda}\right)
\end{equation}
under the output logarithmic constraint
$
\mathbf{Y}\in \mathcal{D}_{{\rm bi}}(A)
$
for some 
prescribed dispersion level $A$, where $A \geq \lambda$. In addition, the corresponding capacity achieving output distribution is
\begin{equation}
\mathbf{Y}^* \sim {\rm Cauchy}_2(0,{\rm diag}(A^2,A^2)).
\end{equation}
\label{thm:T2}
\end{thm}

\begin{proof}
%Next we prove Theorem~\ref{thm:T2}. 
For the additive channel we are considering, the mutual information can be written as
\begin{equation}
I(\mathbf{X};\mathbf{Y})=h(\mathbf{Y})-h(\mathbf{N}).
\end{equation}
The problem of finding the channel capacity turns out to be an optimization problem:
$
C(A,\lambda)=\sup_{\mathbf{Y}\in\mathcal{D}_{\text{bi}}(A)} I(\mathbf{X};\mathbf{Y}).
$
We have
\begin{align}
    C(A,\lambda)
    &=\sup_{\mathbf{Y} \in \mathcal{D}_{\text{bi}}(A)} \Big\{h(\mathbf{Y} ) - h(\mathbf{N} )\Big\}\\
    &=\Big\{\sup_{\mathbf{Y} \in \mathcal{D}_{\text{bi}}(A)} h(\mathbf{Y} )\Big\} - \ln(2\pi e^3\lambda^2), \label{eq:54}
\end{align}
where
\begin{align}
    \sup_{\mathbf{Y} \in \mathcal{D}_{\text{bi}}(A)} h(\mathbf{Y})
    &=
    \sup_{k\in[\lambda,A]}
    \Big\{ \sup_{\mathbf{Y}\in \mathcal{D}_{k,\text{bi}}} h(\mathbf{Y}) \Big\}\label{eq:55}\\
    &=
    \sup_{k\in[\lambda,A]}
    \ln(2\pi e^3 k^2) \label{eq:56}\\
    &= \ln(2\pi e^3 A^2). \label{eq:57}
\end{align}
Combining equations \eqref{eq:54} and \eqref{eq:57} yields
\begin{align}
\begin{split}
C(A,\lambda)
=&\ \ln(2\pi e^3 A^2) - \ln(2\pi e^3 \lambda^2)\\
=&\ \ln\left(\dfrac{A^2}{\lambda^2}\right)
=2\ln\left(\dfrac{A}{\lambda}\right).
\end{split}
\end{align}
Notice that the equalities \eqref{eq:55}-\eqref{eq:57} occur if and only if $\mathbf{Y}$ distributes as ${\rm Cauchy}_2(0,\text{diag}(A^2,A^2))$. Hence, we have proved Theorem~\ref{thm:T2}.
\end{proof}
\color{black}

We encapsulate the primary findings of this section in Table~I. Note that within this context, $A$ represents the output signal dispersion level specific to FAP channels. Moreover, by defining $A^2$ as $\sigma^2+P$, we establish a correspondence between the Gaussian case and the 2D FAP case. To visually illustrate these capacity curves, we present a graphical representation in Fig.~2.

%We summarize the main results of this section into the following Table~I.
%Note that $A$ is the dispersion level for FAP channels.
%Furthermore, if we identify $A^2:=\sigma^2+P$, the Gaussian case will coincide with the 2D FAP case.
%We also provide a graphical demonstration of these capacity curves in Fig.~2.

% Insert your table here
\begin{table}[ht]
\caption{2D and 3D FAP channel capacity curves\\ with a comparison to Gaussian channels}
\label{tab:main}
\centering
\begin{tabular}{|r||c|c|c|}
\hline
\textbf{Channel Type} & \textbf{Gaussian} & \textbf{2D FAP} & \textbf{3D FAP} \\ 
\hline
Constraint & $L^2$ constr. & log constr. & log constr. \\ 
\hline
Capacity Formula & $\frac12\ln\left(\dfrac{\sigma^2+P}{\sigma^2}\right)$ & $\ln\left(\dfrac{A}{\lambda}\right)$ & $2\ln\left(\dfrac{A}{\lambda}\right)$\\ 
\hline
Parameter Range & N/A & $A\geq\lambda$ & $A\geq\lambda$\\
\hline
\end{tabular}
\label{tab:yourtable}
\end{table}

%%--
\begin{figure}[!t]  %[htbp] 
\label{fig:capacity}
\centerline{\includegraphics[width=0.48\textwidth]{./Figures/ycl-icc2024-fig2-large.eps}}
\caption{
Illustration depicting the capacity curves in relation to the output signal dispersion level $A$ for both 2D and 3D FAP channels. It is worth noting that the Gaussian curve coincides with the 2D FAP curve when we equate $A^2$ to $\sigma^2+P$ in the Gaussian case.
}
%\caption{
%Graphical demonstration of the capacity curves versus the output signal dispersion level $A$ for 2D and 3D FAP channels.
%Note that the Gaussian curve and 2D FAP curve coincide if we identify $A^2:=\sigma^2+P$ in the Gaussian case.
%}
\label{}
\end{figure}
%%--

\section{Conclusion}\label{sec:conclude}
\color{black}
This study examines the capacity problem for zero-drift first arrival position (FAP) communication channel in diffusive molecular communication (MC) systems. 
Our findings reveal that when the drift velocity approaches zero, both the 2D and 3D FAP channels reduce to univariate and bivariate Cauchy distributions, respectively. 

Although a recent work \cite{lee2023char} investigate the FAP channel capacity under the case of positive \emph{vertical drift}, at the time of this writing, the determination of the Shannon capacity for the FAP channels with \emph{zero drift} (i.e. the multidimensional Cauchy channels) remains an open problem. This challenge stems from the heavy-tailed nature of the Cauchy distribution, which belongs to a broader family of $\alpha$-stable distributions that often lack a second moment. As a result, traditional energy constraints such as the variance of random variables are inadequate for characterizing the signal space.

To address this challenge, we introduce a relative power measure, the $\alpha$-power, into the field of molecular communication. We explicitly define the signal space required to characterize the zero-drift FAP channel capacity and derive the Shannon capacity for the FAP channel in 2D and 3D spaces in closed-form under a fresh logarithmic constraint. The capacity formulations articulated in Theorem~\ref{thm:T1} and Theorem~\ref{thm:T2} underscore an intriguing phenomenon: the 3D FAP channel's capacity eclipses its 2D counterpart by a factor of two when benchmarked against the identical dispersion parameter \(A\). This revelation amplifies the perspective that as spatial dimension \(n\) expands, FAP channels inherently harness greater informational capacity.
\color{black}


%\input{./section-intro-v8}
%\input{./section2-v8}
%\input{./section3-v8} % without Conclusion
%\input{./section4-v8} % This is the Conclusion

\bibliographystyle{IEEEtran}
% Your bibliography goes here
\balance
\bibliography{main}

\appendices
\section{Density Function and Entropy of Multivariate Cauchy Distribution}
\label{appendix-A}
The Cauchy distribution, named after Augustin Cauchy, is a continuous probability distribution often used in statistics as a canonical example of distribution since both its expected value and its variance are undefined. It is also known, especially among physicists, as the Lorentz distribution.
In mathematics, it is closely related to the fundamental solution for the Laplace equation in the half-plane, and it is one of the few distributions that is stable (see \cite{pierce1997application,fahs2012capacity}) and has a density function that can be expressed analytically. (Other examples are normal distribution and L\'{e}vy distribution.)

The (univariate) Cauchy distribution has the probability density function (PDF) which can be expressed as:
\begin{equation}
    f(x;x_0,\gamma) =
    \dfrac{1}{\pi \gamma} 
    \dfrac{\gamma^2}
    {(x-x_0)^2+\gamma^2
    }
    =
    \dfrac{1}{\pi \gamma} 
    \dfrac{1}
    {1+(\frac{x-x_0}{\gamma})^2
    },
\label{eq:ch-dist}
\end{equation}
where $x_0\in \mathbb{R}$ is the location parameter, and $\gamma>0$ is the scaling parameter (see \cite{shao1993signal}). 
For the purpose of exploring the channel capacity, since the location parameter $x_0$ is irrelevant to the entropy, we may assume without loss of generality that $x_0=0$, yielding
the so called \textit{symmetry Cauchy} with PDF as:
\begin{equation}
    f(x;\gamma) =
    \dfrac{\gamma}{\pi} 
    \dfrac{1}
    {x^2+\gamma^2
    }
    =
    \dfrac{1}{\pi \gamma} 
    \dfrac{1}
    {1+(\frac{x}{\gamma})^2
    }.
\label{eq:sm-ch-dist}
\end{equation}
The entropy of Cauchy distribution can be evaluated by direct calculation. From Eq.~\eqref{eq:ch-dist}, we have
\begin{equation}
    h(X)=\ln(4\pi \gamma)
\end{equation}
whenever $X\sim \rm{Cauchy}(x_0,\gamma)$, and $x_0$ can be chosen arbitrarily.


As for multivariate Cauchy distribution, our notation system mainly follows \cite{kotz2004multivariate}.
Consider a $p$-dimensional Euclidean random vector $\mathbf{X}=(X_1,\cdots,X_p)^\tsp$ which follows multivariate Cauchy distribution.
We use the notation $\mathbf{X} \sim {\rm Cauchy}_p(\boldsymbol\mu,\Sigma)$ to specify the parameters, where $\boldsymbol\mu\in\R^{p\times 1}$ is the location vector, and $\Sigma\in\R^{p\times p}$ is the scale matrix describing the shape of the distribution.
The PDF of multivariate Cauchy distribution is given by
the following formula:
\begin{align} 
\label{eq:pvariate-ch-dist}
\begin{split}
    &f_{\mathbf{X}}^{(p)}(\mathbf{x};\boldsymbol\mu,\Sigma)\\
    =&\ 
    \dfrac{\Gamma(\frac{1+p}{2})}
    {\Gamma(\frac{1}{2})\pi^{\frac{p}{2}}
    |\Sigma|^{\frac{1}{2}}
    [1+(\mathbf{x}-\boldsymbol\mu)^\top\Sigma^{-1}(\mathbf{x}-\boldsymbol\mu)]^\frac{1+p}{2}
    }.
\end{split}
\end{align}
Note that $\Sigma$ is by nature a positive-definite square matrix.
For our later purpose, we mainly work in the case $p=2$. We provide the PDF of this special case here for convenience:
\begin{equation}
    f_{\mathbf{X}}^{(2)}(\mathbf{x};\boldsymbol\mu,\Sigma)=
    \dfrac{\Gamma(\frac{3}{2})}
    {\pi^{\frac{3}{2}}
    |\Sigma|^{\frac{1}{2}}
    [1+(\mathbf{x}-\boldsymbol\mu)^\top\Sigma^{-1}(\mathbf{x}-\boldsymbol\mu)]^\frac{3}{2}
    }.
\label{eq:bivariate-ch-dist}
\end{equation}
Eq.~\eqref{eq:bivariate-ch-dist} is the so-called bivariate Cauchy distribution.

Similar to the univariate case,
we can without loss of generality set $\boldsymbol\mu=\mathbf{0}$ in equations \eqref{eq:pvariate-ch-dist} and \eqref{eq:bivariate-ch-dist} for the purpose of entropy and channel capacity analysis.
When $\boldsymbol\mu = \mathbf{0}$, the multivariate Cauchy is called \textit{central}. For the case that the Cauchy is central, we abuse the notations ${\rm R}$ and $\Sigma$ to mean the same scale matrix.
We shall assume all Cauchy distributions considered in this paper to be central from now on. In addition, for the special case that $\Sigma$ is a diagonal matrix, 
namely
\begin{equation}
\Sigma=\text{diag}(\gamma^2,\gamma^2)
=\begin{bmatrix}
\gamma^2&0\\
0&\gamma^2
\end{bmatrix},
\end{equation}
the central bivariate Cauchy density function can be simplified into a more concise form:
\begin{equation}
\dfrac{1}{2\pi}
\dfrac{\gamma}{
(n_1^2+n_2^2+\gamma^2)^{3/2}
},
\label{eq:simple-bi-ch}
\end{equation}
where $\mathbf{n}=(n_1,n_2)$ is the Cauchy random vector.

Finally, the (differential) entropy of $p$-variate central Cauchy distribution can be found in [4]. We briefly state the results here for later use. Suppose the scale matrix ${\rm R}$ of a $p$-variate Cauchy $\mathbf{X}$ is prescribed, the entropy of $\mathbf{X}$ can be expressed as
\begin{equation}
    h(\mathbf{X};{\rm R})=\dfrac{1}{2}\ln|{\rm R}|+\Phi(p),
\end{equation}
where $|{\rm R}|$ represents the determinant of matrix ${\rm R}$; $\Phi(p)$ is a constant depending only on dimension $p$, and it can be evaluated through
\begin{align}
\begin{split}
    \Phi(p)
    =&\ \ln
    \left[
    \dfrac{\pi^{\frac{p}{2}}}{\Gamma(\frac{p}{2})}B\left(\frac{p}{2},\dfrac{1}{2}\right)
    \right]\\
    &\quad +\dfrac{1+p}{2}\left[
    \psi\left(\frac{1+p}{2}\right)-\psi\left(\frac12\right)
    \right].
\end{split}
\end{align}
In the above expression, $B(x,y)=\frac{\Gamma(x)\Gamma(y)}{\Gamma(x+y)}$ is the so-called \textit{beta function}, and
\begin{equation}
\psi(t):=\dfrac{d}{dt}\big[\ln\Gamma(t)\big]   
\end{equation}
is known as the \textit{digamma function}.


\section{Proof of Theorem~\ref{thm:T1}}
In this appendix section, we prove Theorem~\ref{thm:T1} step by step.
For the additive channel
\begin{equation}
    Y=X+N
\end{equation}
we are considering, the mutual information can be written as
$I(X;Y)=h(Y)-h(N)$ according to Section~\ref{pre-A}. The issue of finding the capacity of channel \eqref{eq:29} turns out to be an optimization problem:
\begin{align}
    C(A,\lambda)&=\sup_{Y\in\mathcal{D}(A)} I(X;Y).
\end{align}
The calculation is as follows. We have
\begin{align}
    C(A,\lambda)
    &=\sup_{Y \in \mathcal{D}(A)} \Big\{h(Y) - h(N)\Big\}\\
    &=\Big\{\sup_{Y \in \mathcal{D}(A)} h(Y)\Big\} - \ln(4\pi\lambda), \label{eq:39}
\end{align}
where
\begin{align}
    \sup_{Y \in \mathcal{D}(A)} h(Y)
    &=
    \sup_{k\in[\lambda,A]}
    \Big\{ \sup_{Y\in \mathcal{D}_k} h(Y) \Big\}\label{eq:40}\\
    &=
    \sup_{k\in[\lambda,A]}
    \ln(4\pi k) \label{eq:41}\\
    &= \ln(4\pi A). \label{eq:42}
\end{align}
Combining equations \eqref{eq:39} and \eqref{eq:42} yields
\begin{equation}
    C(A,\lambda) = 
    \ln(4\pi A) - \ln(4\pi \lambda) = \ln\left(\dfrac{A}{\lambda}\right).
\end{equation}
Notice that the equalities \eqref{eq:40}-\eqref{eq:42} hold if and only if $Y$ distributes as ${\rm Cauchy}(0,A)$. Hence, Theorem~\ref{thm:T1} is proved.
\label{appendix-B}

%% This is actually Appendix-C
\section{Some Useful Facts about Cauchy Distribution}\label{pre-B}
\color{black}
In this appendix, we delve into specific properties of Cauchy distributions, building upon the foundational probability density function formula and entropy evaluations presented in Appendix~\ref{appendix-A}. 

The first two properties initially highlight the behavior of the independent sum of Cauchy distributions. To elucidate, the univariate symmetric Cauchy distribution remains invariant under independent addition, as articulated in the subsequent lemma.

\begin{lma}
Given two independent Cauchy random variables, $U \sim {\rm Cauchy}(0,\sigma)$ and $V \sim {\rm Cauchy}(0,\tau)$, their sum is described as:
\begin{equation}
Z = U + V \sim {\rm Cauchy}(0,\sigma+\tau).    
\end{equation}
\label{lma:L1}
\end{lma}
\noindent When considering the independent summation of bivariate Cauchy distributions, the lemma below encapsulates the relationship.

\begin{lma}
With $\mathbf{U} \sim {\rm Cauchy}_2\Big(\mathbf{0},\Sigma_1={\rm diag}(\sigma^2,\sigma^2)\Big)$ and $\mathbf{V} \sim {\rm Cauchy}_2\Big(\mathbf{0},\Sigma_2={\rm diag}(\tau^2,\tau^2)\Big)$ representing two independent bivariate Cauchy random vectors, their summation follows:
\begin{align}
\begin{split}
\mathbf{Z} &= \mathbf{U} + \mathbf{V} \\
&\sim {\rm Cauchy}_2\Big(\mathbf{0},\Sigma_3={\rm diag}\left((\sigma+\tau)^2,(\sigma+\tau)^2\right)\Big).
\end{split}
\end{align}
\label{lma:L2}
\end{lma}
\noindent Note that in Lemma~2, we employ the notation
\begin{equation}
{\rm diag}(a,b) =
\begin{bmatrix}
a & 0 \\
0 & b
\end{bmatrix}
\end{equation}
to succinctly represent diagonal matrices.

The final property examines the linear combination of components of a Cauchy vector, \(\mathbf{X}\). Given an arbitrary constant vector \(\mathbf{v} \in \R^{p \times 1}\), the lemma below delineates the relationship.

\begin{lma}
For a vector $\mathbf{X} \sim {\rm Cauchy}_p(\boldsymbol\mu,\Sigma)$, the relation is given by:
\begin{equation}
\mathbf{v}^\intercal \mathbf{X}
\sim {\rm Cauchy}(\mathbf{v}^\intercal \boldsymbol\mu,
\mathbf{v}^\intercal \Sigma \mathbf{v}),
\end{equation}
where \(\boldsymbol{\mu}\) signifies the location vector of \(\mathbf{X}\) and \(\Sigma\) represents its scale matrix.
\label{lma:L3}
\end{lma}
\color{black}

%%--------------
\begin{comment}
For the probability density function formula and entropy evaluation of Cauchy distribution, please refer to Appendix~\ref{appendix-A}. In this appendix section, we enumerate three useful properties of Cauchy distributions.

The first and second property is about the independent sum of two Cauchy distributions. Briefly speaking, univariate symmetry Cauchy distribution is closed under independent sum, as the following lemma states.
\begin{lma}
Letting $U \sim {\rm Cauchy}(0,\sigma)$,
$V \sim {\rm Cauchy}(0,\tau)$ be two independent Cauchy random variables, we have
\begin{equation}
Z=U+V \sim {\rm Cauchy}(0,\sigma+\tau).    
\end{equation}
\label{lma:L1}
\end{lma}
\noindent As for the independent sum property of bivariate Cauchy distributions, we have the following lemma.
\begin{lma}
Letting $\mathbf{U} \sim {\rm Cauchy}_2\Big(\mathbf{0},\Sigma_1={\rm diag}(\sigma^2,\sigma^2)\Big)$,
$\mathbf{V} \sim {\rm Cauchy}_2\Big(\mathbf{0},\Sigma_2={\rm diag}(\tau^2,\tau^2)\Big)$ be two independent bivariate Cauchy random vectors, we have
\begin{align}
\begin{split}
\label{lma:L2}
\mathbf{Z} &=\mathbf{U}+\mathbf{V} \\
&\sim {\rm Cauchy}_2\Big(\mathbf{0},\Sigma_3={\rm diag}\left((\sigma+\tau)^2,(\sigma+\tau)^2\right)\Big).
\end{split}
\end{align}
\end{lma}
\noindent Notice that in Lemma~2, we have used the notation
\begin{equation}
{\rm diag}(a,b)=
\begin{bmatrix}
a & 0\\
0 & b
\end{bmatrix}
\end{equation}
for conveniently representing diagonal matrices.

The third property is about
linear combination of components of Cauchy vector $\mathbf{X}$. Let $\mathbf{v}\in\R^{p\times 1}$ be an arbitrary constant vector, then the following lemma holds.
\begin{lma}
If $\mathbf{X} \sim {\rm Cauchy}_p(\boldsymbol\mu,\Sigma)$, then
\begin{equation}
\mathbf{v}^\intercal \mathbf{X}
\sim {\rm Cauchy}(\mathbf{v}^\intercal \boldsymbol\mu,
\mathbf{v}^\intercal \Sigma \mathbf{v}),
\end{equation}
where $\boldsymbol{\mu}$ is the location vector of $\mathbf{X}$ and $\Sigma$ is the scale matrix of $\mathbf{X}$.
\label{lma:L3}
\end{lma}
\end{comment}

\end{document}
